% from ega0-20-fr.tex
\oldpage[0\textsubscript{IV-1}]{121}
\begin{corollary}[20.2.3]
\label{0.20.2.3-fr}
Soient $A$, $B$, $C$ trois anneaux commutatifs, $u:A\to B$, $v:B\to C$
deux homomorphismes d'anneaux, $L$ un $C$-module. On a une suite exacte
canonique de $A$-modules
\begin{align}
\tag{20.2.3.1}
\label{0.20.2.3.1-fr}
  0\to\operatorname{D\acute{e}r}_B(C,L)\xrightarrow{u^0}
  \operatorname{D\acute{e}r}_A(C,L)&\xrightarrow{v^0}
  \operatorname{D\acute{e}r}_A(B,L)\xrightarrow{\partial} \\
  &\to\operatorname{Exalcom}_B(C,L)\xrightarrow{u^1}
  \operatorname{Exalcom}_A(C,L)\xrightarrow{v^1}
  \operatorname{Exalcom}_A(B,L) \notag
\end{align}
fonctorielle en $L$.
\end{corollary}

\begin{proof}
Le raisonnement est le même que dans \sref{0.20.2.2-fr} une fois qu'on a
vérifié que pour toute dérivation $D\in\operatorname{D\acute{e}r}_A(B,L)$,
$\partial(D)$ est bien la classe d'une $B$-extension de $C$ par $L$ qui est
une $B$-algèbre commutative ; mais cela résulte aussitôt de la commutativité
de $C$ et du fait que $L$ est un $C$-module.
\end{proof}

\begin{corollary}[20.2.4]
\label{0.20.2.4-fr}
Sous les hypothèses de \sref{0.20.2.2-fr} (resp.
\sref{0.20.2.3-fr}), on a une suite exacte canonique, fonctorielle en $L$,
\begin{equation}
\tag{20.2.4.1}
\label{0.20.2.4.1-fr}
  0\to\operatorname{D\acute{e}r}_B(C,L)\xrightarrow{u^0}
  \operatorname{D\acute{e}r}_A(C,L)\xrightarrow{v^0}
  \operatorname{D\acute{e}r}_A(B,L)\xrightarrow{\partial}
  \operatorname{Exan}_{B/A}(C,L)\to0
\end{equation}
(resp.
\begin{equation}
\tag{20.2.4.2}
\label{0.20.2.4.2-fr}
  0\to\operatorname{D\acute{e}r}_B(C,L)\xrightarrow{u^0}
  \operatorname{D\acute{e}r}_A(C,L)\xrightarrow{v^0}
  \operatorname{D\acute{e}r}_A(B,L)\xrightarrow{\partial}
  \operatorname{Exalcom}_{B/A}(C,L)\to0\text{).}
\end{equation}
Cela résulte de la définition de $\operatorname{Exan}_{B/A}(C,L)$ (resp.
$\operatorname{Exalcom}_{B/A}(C,L)$) ((18.3.7) et (18.4.2)).
\end{corollary}

\begin{remark}[20.2.5]
\label{0.20.2.5-fr}
Supposons que l'on ait un diagramme commutatif d'homomorphismes d'anneaux
\[
\xymatrix{
  A \ar[r] & B \ar[r] & C \\
  A' \ar[u] \ar[r] & B' \ar[u] \ar[r] & C' \ar[u]
}
\]
Alors on a un diagramme commutatif
\[
\xymatrix@C=0.5em{
  0 \ar[r] & \operatorname{D\acute{e}r}_B(C,L) \ar[r] \ar[d] &
  \operatorname{D\acute{e}r}_A(C,L) \ar[r] \ar[d] &
  \operatorname{D\acute{e}r}_A(B,L) \ar[r]^-\partial \ar[d] &
  \operatorname{Exan}_B(C,L) \ar[r] \ar[d] &
  \operatorname{Exan}_A(C,L) \ar[r] \ar[d] &
  \operatorname{Exan}_A(B,L) \ar[d] \\
  0 \ar[r] & \operatorname{D\acute{e}r}_{B'}(C',L) \ar[r] &
  \operatorname{D\acute{e}r}_{A'}(C',L) \ar[r] &
  \operatorname{D\acute{e}r}_{A'}(B',L) \ar[r]^-\partial &
  \operatorname{Exan}_{B'}(C',L) \ar[r] &
  \operatorname{Exan}_{A'}(C',L) \ar[r] &
  \operatorname{Exan}_{A'}(B',L)
}
\]
et de même pour les suites exactes \sref{0.20.2.3.1-fr},
\sref{0.20.2.4.1-fr} et \sref{0.20.2.4.2-fr}.
\end{remark}

\subsection{Dérivations continues dans les anneaux topologiques}
\label{subsection:0.20.3-fr}

\begin{env}[20.3.1]
\label{0.20.3.1-fr}
Étant donnés deux anneaux topologiques $A$, $B$ (linéairement topologisés
comme toujours), nous désignerons par $\operatorname{Hom.\ cont}(A,B)$
l'ensemble des homomorphismes continus de $A$ dans $B$. Étant donné un anneau
topologique $A$, la catégorie des $A$-anneaux topologiques se définit comme
celle des $A$-anneaux (18.1.4) en remplaçant partout « anneau » par « anneau
topologique » et « homomorphisme » par « homomorphisme continu » ; si $B$

